% !TeX root = ../../Book2.tex
% 纲要标题：### 8.3 对称代数

\section{对称代数}
\label{b2:sec:0803}

张量代数保留因子的排列顺序。若把 \(uv\) 与 \(vu\) 识别，就得到只记录交换乘积的代数。这个商构造在任意交换底环上都成立；把它认成张量幂中的“不变部分”，却通常需要能除以阶乘。本节把这两个对象分别定义，再说明何时存在指定的同构。

% 纲要要点（供目录核对）：
%
% * 对称代数的泛性质
% * 有限自由模上的多项式解释
% * 对称幂的余不变量解释、对称化及特征 p 的取幂
% * 三次对称化的轨道计算及特征二、三的比较
% * 张量代数与对称代数的扩张标量相容性
%

\subsection{对称代数的泛性质}
\label{b2:subsec:0803:01}

\begin{definition}\label{b2:def:symmetric-algebra}
设 \(R\) 为交换含幺环，\(M\) 为幺性 \(R\) 模。令 \(J\) 为 \(T_R(M)\) 中由全部
\[
u\otimes v-v\otimes u,\qquad u,v\in M,
\]
生成的双边理想。商代数
\[
\operatorname{Sym}_R(M)\coloneqq T_R(M)/J
\]
称为 \(M\) 的\term[duichen-daishu]{对称代数}[symmetric algebra]。由齐次二次关系及\cref{b2:prop:homogeneous-ideal-quotient}，它有自然分次；对每个 \(d\ge0\)，其次数 \(d\) 的分量记为 \(\operatorname{Sym}_R^d(M)\)，称为 \(M\) 的 \(d\) 次\term[duichen-mi]{对称幂}[symmetric power]。张量 \(m_1\otimes\cdots\otimes m_d\) 的像记为 \(m_1\cdots m_d\)。
\end{definition}
\symidx{\(\operatorname{Sym}_R(M)\)：模的对称代数}
\symidx{\(\operatorname{Sym}_R^d(M)\)：模的 \(d\) 次对称幂}

关系在次数二，故次数零与一次部分仍分别为 \(R,M\)。全部一次生成元在商中彼此交换；任意两个字可逐个交换其中的生成元，故整个商代数交换。任何从 \(T_R(M)\) 到交换 \(R\) 代数的同态都将这些交换关系送到零，因而唯一经过这个商。所以 \(\operatorname{Sym}_R(M)\) 是强制一次生成元彼此交换得到的普遍交换商。一般模中的线性关系仍予保留，并没有因交换乘法而消失。

\begin{theorem}[对称代数的泛性质]\label{b2:thm:symmetric-algebra-universal}
设 \(R\) 为交换含幺环，\(M\) 为幺性 \(R\) 模。对任意交换含幺 \(R\) 代数 \(A\) 及 \(R\) 线性映射 \(f:M\rightarrow A\)，存在唯一保幺 \(R\) 代数同态
\[
\widehat f:\operatorname{Sym}_R(M)\longrightarrow A
\]
在一次分量上等于 \(f\)。它将 \(m_1\cdots m_d\) 送到 \(f(m_1)\cdots f(m_d)\)。
\end{theorem}
\begin{proof}
\cref{b2:thm:tensor-algebra-universal}先把 \(f\) 延拓为 \(T_R(M)\rightarrow A\)。每个关系的像为 \(f(u)f(v)-f(v)f(u)=0\)，因为 \(A\) 交换。因此延拓同态将 \(J\) 送到零，诱导唯一商同态。反过来，商中的一次分量和结构标量生成全代数，所需同态在它们上的值已指定，故唯一。
\end{proof}

对线性映射 \(f:M\rightarrow N\)，上述定理给出分次同态 \(\operatorname{Sym}_R(f)\)，其在 \(d\) 次部分的公式为
\[
\operatorname{Sym}_R^d(f)(m_1\cdots m_d)=f(m_1)\cdots f(m_d).
\]
这些映射保持恒等与复合，因为在一次生成元上保持。因此对称代数与每次对称幂都是函子。目标交换的条件不能删去：若 \(M=R^2\)，把两个基向量任意映到非交换代数中的两个元素，只有它们交换时才可经 \(\operatorname{Sym}_R(M)\) 延拓。

\ProseParagraphBreak
在次数 \(d\) 中，乘积 \(m_1\cdots m_d\) 不再记录各因子的排列。因此要把张量与它的置换像识别，也就是商去它们的差；这不同于从张量幂中选出已经被每个置换固定的元素。

\begin{proposition}\label{b2:prop:symmetric-power-universal}
设 \(R\) 为交换含幺环，\(M,N\) 为幺性 \(R\) 模，\(d\ge1\)。令 \(S_d\) 通过置换因子作用于 \(M^{\otimes d}\)，\(q_d:M^{\otimes d}\rightarrow\operatorname{Sym}_R^d(M)\) 为取商映射，则 \(q_d\) 的核由所有 \(t-P_\sigma t\)（\(t\in M^{\otimes d}\)、\(\sigma\in S_d\)）生成。因此 \(q_d\) 给出自然同构
\[
\operatorname{Sym}_R^d(M)\cong
M^{\otimes d}/\langle t-P_\sigma t\mid t\in M^{\otimes d},\ \sigma\in S_d\rangle_R.
\]
从 \(\operatorname{Sym}_R^d(M)\) 到 \(N\) 的 \(R\) 线性映射，与满足
\[
b(m_{\sigma(1)},\ldots,m_{\sigma(d)})=b(m_1,\ldots,m_d)
\quad(\sigma\in S_d,\ m_i\in M)
\]
的多线性映射 \(b:M^d\rightarrow N\) 一一对应。这样的 \(b\) 称为\term[duichen-duoxianxing-yingshe]{对称多线性映射}[symmetric multilinear map]。
\end{proposition}
\begin{proof}
设 \(K_d\) 为所述差生成的子模。先看相邻两个位置的交换：在纯张量上，\(t-P_{(i\ i+1)}t\) 是一个二次交换关系左右各乘其余纯张量的结果，因此在 \(J\) 中。任意置换均可通过相邻交换完成。若这些交换依次给出纯张量 \(t=t_0,t_1,\ldots,t_r=P_\sigma t\)，则每个 \(t_{j-1}-t_j\) 都属于 \(J\)，而 \(t-P_\sigma t=\sum_{j=1}^r(t_{j-1}-t_j)\)，故它也属于 \(J\)。这先对纯张量成立，再由线性对所有 \(t\) 成立，故 \(K_d\subseteq J\cap M^{\otimes d}\)。

反向，将 \(J\) 的元素写成 \(a(u\otimes v-v\otimes u)b\) 的有限和，再取次数 \(d\) 的分量，并把 \(a,b\) 展开成纯张量的有限和。所得每项都是某个纯张量与相邻交换后的张量之差，所以属于 \(K_d\)。因此核恰为 \(K_d\)。对任意 \(R\) 线性映射 \(f:M\to M'\)，有 \(f^{\otimes d}P_\sigma=P_\sigma f^{\otimes d}\)，故该商同构与 \(\operatorname{Sym}_R^d(f)\) 相容，确为自然同构。

多线性映射先由\cref{b2:thm:multiple-tensor-universal}唯一合成 \(M^{\otimes d}\rightarrow N\)。变量置换不改变值，恰好要求它在上述核上为零，因而唯一经 \(q_d\) 分解。反向复合显然逐变量线性且对称。
\end{proof}

\subsection{有限自由模的多项式解释}
\label{b2:subsec:0803:02}

\begin{theorem}[多项式解释]\label{b2:thm:symmetric-polynomial-algebra}
设 \(R\) 为交换含幺环，\(M\) 为以 \(e_1,\ldots,e_r\) 为基的幺性有限自由 \(R\) 模，则存在唯一分次 \(R\) 代数同构
\[
\operatorname{Sym}_R(M)\longrightarrow R[X_1,\ldots,X_r],
\qquad e_i\longmapsto X_i.
\]
\end{theorem}
\begin{proof}
将基向量送到不定元给出线性映射 \(M\rightarrow R[X_1,\ldots,X_r]\)。目标交换，故\cref{b2:thm:symmetric-algebra-universal}给出正向同态。反向把多项式 \(\sum_\alpha a_\alpha X^\alpha\) 送到 \(\sum_\alpha a_\alpha e_1^{\alpha_1}\cdots e_r^{\alpha_r}\)。和为有限和，商中的 \(e_i\) 彼此交换，因此多项式的加法与乘法公式保证它是代数同态。两复合分别在各 \(e_i\) 与 \(X_i\) 上为恒等，并在结构标量上为恒等；生成性使两复合在全代数为恒等。两映射保持生成元的次数，故保持分次。
\end{proof}

\begin{proposition}\label{b2:prop:symmetric-power-basis}
设 \(R\) 为交换含幺环，\(M\) 为以 \(e_1,\ldots,e_r\) 为基的幺性有限自由 \(R\) 模，\(d\ge0\)，则次数 \(d\) 的单项式
\[
e_1^{a_1}\cdots e_r^{a_r},\qquad a_i\ge0,\quad a_1+\cdots+a_r=d,
\]
构成 \(\operatorname{Sym}_R^d(M)\) 的一组基。若 \(R\ne0\)、\(r\ge1\)，则其秩为 \(\binom{r+d-1}{d}\)。当 \(r=0\) 时，零次部分为 \(R\)，所有正次部分为零。
\end{proposition}
\begin{proof}
多项式作为有限形式和，其不同单项式的系数唯一；上一同构把全次数 \(d\) 的单项式送到题述元素，所以它们构成基。计算数目时，把 \(d\) 个相同的记号排成一行，并插入 \(r-1\) 道隔板。第 \(i\) 段记号的个数就是 \(a_i\)，允许空段。这与从 \(r+d-1\) 个位置选出 \(d\) 个记号位置一一对应，故数目为所述二项式系数。零模情形直接由定义得到。
\end{proof}

例如 \(R^2\) 的三次对称幂以 \(e_1^3,e_1^2e_2,e_1e_2^2,e_2^3\) 为基，秩为四；相应三次张量幂则有八个有序纯基张量。对称商把 \(e_1\otimes e_1\otimes e_2\) 的三个不同排列都映到 \(e_1^2e_2\)；只有把这三个张量先相加，其像才是 \(3e_1^2e_2\)。单个张量的像不需要乘上排列数。

\ProseParagraphBreak
对有限维 \(k\) 向量空间 \(V\)，\(\operatorname{Sym}_k(V)\) 的一次生成元是向量，而 \(V\) 上的线性坐标函数属于对偶空间 \(V^*\)。因此，由线性坐标函数生成的形式多项式代数应写为 \(\operatorname{Sym}_k(V^*)\)：线性泛函 \(\lambda\) 在 \(v\) 处取值为 \(\lambda(v)\)，乘积 \(\lambda_1\cdots\lambda_d\) 则取值为 \(\prod_j\lambda_j(v)\)。选定 \(V\) 的基后，其对偶基对应熟悉的坐标不定元。两空间虽同维，\(V\) 与 \(V^*\) 却没有由向量空间本身指定的识别，不能把向量生成元直接当作坐标函数。

\ProseParagraphBreak
形式多项式也不等同于它定义的函数。例如在 \(\mathbb F_p\) 上，\(X^p-X\) 是非零多项式，却在每个域元素处取零。对称代数采用形式多项式的乘法和系数，不因所有取值为零就把它的元素置零；\secref{b2:sec:0804}将说明这一区别如何影响不变环的定义。

\subsection{对称张量、对称商与对称化}
\label{b2:subsec:0803:03}

先考虑一个交换两个基向量的作用：在自由模 \(Ru\oplus Rv\) 中，不变子模为 \(R(u+v)\)，余不变量则为 \((Ru\oplus Rv)/R(u-v)\)。前者要求两个坐标本来相同，后者让 \(u,v\) 的类相同。下面把这两种构造用于张量因子的置换。

\ProseParagraphBreak
设 \(R\) 为交换含幺环，\(M\) 为幺性 \(R\) 模，\(d\ge1\)。令 \(S_d\) 按\cref{b2:prop:multiple-tensor-functoriality}的公式作用在 \(M^{\otimes d}\) 上。这一作用的\term[bubianliang]{不变量}[invariants]是子模
\[
(M^{\otimes d})^{S_d}
=\{\,t\in M^{\otimes d}\mid P_\sigma t=t\;\text{对所有}\;\sigma\in S_d\,\},
\]
其中的元素称为\term[duichen-zhangliang]{对称张量}[symmetric tensor]。同一作用的\term[yububianliang]{余不变量}[coinvariants]是商模
\[
(M^{\otimes d})_{S_d}
\coloneqq M^{\otimes d}/\langle t-P_\sigma t\mid t\in M^{\otimes d},\ \sigma\in S_d\rangle_R.
\]
由\cref{b2:prop:symmetric-power-universal}，商映射给出 \(\operatorname{Sym}_R^d(M)\cong(M^{\otimes d})_{S_d}\)。因此包含与取商组成
\[
(M^{\otimes d})^{S_d}\hookrightarrow M^{\otimes d}
\xrightarrow{q_d}\operatorname{Sym}_R^d(M)
\cong(M^{\otimes d})_{S_d},
\]
其中 \(q_d\) 满射，总有指定映射
\[
q_d\big|_{(M^{\otimes d})^{S_d}}:
(M^{\otimes d})^{S_d}\longrightarrow\operatorname{Sym}_R^d(M).
\]
线性映射的张量幂保持不变子模，并与 \(q_d\) 相容，所以这个限制映射是自然的。问题在于它何时成为同构，而不是两个模是否偶然具有相同的秩或维数。
\symidx{\((M^{\otimes d})^{S_d},(M^{\otimes d})_{S_d}\)：对称张量的不变量与余不变量}

\ProseParagraphBreak
在次数二，若 \(2\) 可逆，令 \(P\) 为交换两个因子的算子，则 \((t+Pt)/2\) 被 \(P\) 固定；若 \(t\) 本来对称，这个平均仍为 \(t\)。一般次数下，对全部 \(d!\) 个置换同样取平均，就能得到一个固定不变张量的投影。

\begin{proposition}\label{b2:prop:symmetrization-characteristic}
设 \(R\) 为交换含幺环，\(M\) 为幺性 \(R\) 模，\(d\ge1\)。若 \(d!\) 在 \(R\) 中可逆，则\term[duichen-hua]{对称化}[symmetrization]算子
\[
E_d=\frac1{d!}\sum_{\sigma\in S_d}P_\sigma
\]
是投影到 \((M^{\otimes d})^{S_d}\) 的线性映射，并诱导自然同构
\[
\overline E_d:\operatorname{Sym}_R^d(M)\xrightarrow{\ \cong\ }(M^{\otimes d})^{S_d}.
\]
其逆是商映射 \(q_d\) 在不变子模上的限制。
\end{proposition}
\begin{proof}
左乘一个固定置换只重新排列求和项，因此 \(P_\pi E_d=E_d\)，像都不变；若 \(t\) 已不变，则 \(E_dt=t\)。所以像正是不变子模，且 \(E_d^2=E_d\)。

同样，\(E_dP_\pi=E_d\)，故 \(E_d\) 零化全部 \(t-P_\pi t\)。由\cref{b2:prop:symmetric-power-universal}，它唯一经 \(q_d\) 分解为 \(\overline E_d\)。在对称商中 \(q_dP_\sigma=q_d\)，因而 \(q_dE_d=q_d\)，说明 \(q_d\overline E_d\) 为恒等；在不变子模上 \(E_d\) 为恒等，又说明 \(\overline E_dq_d\) 为恒等。线性映射的张量幂与每个 \(P_\sigma\) 交换，所以也与平均算子相容；所得同构因此自然。
\end{proof}

阶乘可逆保证上述同构，但它不可逆时须另行检查，不可据此一概断言同构失败。特征零也不保证阶乘可逆：取 \(R=\mathbb Z\)、\(M=\mathbb Ze\oplus\mathbb Zf\)。交换两个因子的不变张量恰为
\[
a(e\otimes e)+b(e\otimes f+f\otimes e)+c(f\otimes f),
\qquad a,b,c\in\mathbb Z,
\]
因为在纯基张量的展开中，\(e\otimes f\) 与 \(f\otimes e\) 的系数必须相同。它在 \(\operatorname{Sym}_{\mathbb Z}^2(M)\) 中的像为
\[
ae^2+2bef+cf^2.
\]
\(e^2,ef,f^2\) 是这个对称幂的一组基，所以 \(ef\) 不在像中，否则需要整数 \(b\) 满足 \(2b=1\)。这里商映射在不变子模上的限制不是满射，障碍是 \(2\) 在 \(\mathbb Z\) 中不可逆。

\ProseParagraphBreak
正特征下还可能出现单射性失败。取 \(k\) 特征为素数 \(p\)，\(V=ke\oplus kf\)，在 \(V^{\otimes p}\) 中令
\[
s=\sum_{j=1}^p e\otimes\cdots\otimes
\underset{\text{第 }j\text{ 个位置}}{f}
\otimes\cdots\otimes e.
\]
各项是互不相同的纯基张量，所以 \(s\ne0\)；置换只重排这些项，所以 \(s\) 不变。但
\[
q_p(s)=p\,e^{p-1}f=0.
\]
因此商映射在不变张量上不是单射。最小例子是特征二下的 \(e\otimes f+f\otimes e\)：它非零且对称，却在二次对称幂中为零。反过来，\(ef\) 仍是对称幂的一个非零基元素。

\begin{example}\label{b2:exm:cubic-symmetrization-orbits}
设 \(k\) 是域，\(V=ke\oplus kf\)。在 \(V^{\otimes3}\) 的八个纯基张量上，\(S_3\) 的置换作用有四个轨道，按 \(f\) 出现零、一、二、三次分类。相应轨道和为
\[
\begin{aligned}
s_0&=e\otimes e\otimes e,\\
s_1&=e\otimes e\otimes f+e\otimes f\otimes e+f\otimes e\otimes e,\\
s_2&=e\otimes f\otimes f+f\otimes e\otimes f+f\otimes f\otimes e,\\
s_3&=f\otimes f\otimes f.
\end{aligned}
\]
置换作用只重排纯基张量，因而一个张量不变，当且仅当同一轨道上的系数相同。四个轨道的支集互不相交，所以 \(s_0,s_1,s_2,s_3\) 线性无关，并生成全部不变张量。它们是一组基；\(\operatorname{Sym}_k^3(V)\) 则以 \(e^3,e^2f,ef^2,f^3\) 为基。限制商映射把 \(s_0,s_1,s_2,s_3\) 分别送到 \(e^3,3e^2f,3ef^2,f^3\)，因此相对于这两组基的矩阵为
\[
\operatorname{diag}(1,3,3,1).
\]
若 \(k=\mathbb Q\)，它是同构，且 \(e^2f\) 的逆像为 \(s_1/3\)。这也等于六个置换的平均：\(e\otimes e\otimes f\) 的每个不同排列出现两次，所以平均为 \(2s_1/6=s_1/3\)。

若 \(\operatorname{char}k=2\)，上述矩阵为单位矩阵，限制映射仍是同构，尽管 \(3!=0\) 不可逆。若 \(\operatorname{char}k=3\)，它的核为 \(ks_1\oplus ks_2\)，像为 \(ke^3\oplus kf^3\)，秩为二；两端都是四维，也不保证指定映射可逆。因此阶乘可逆是统一保证同构的充分条件，并不是一个特定模上限制映射可逆的必要条件。
\end{example}

删去平均公式中的分母，就得到始终有定义的求和算子 \(N_d=\sum_\sigma P_\sigma\)。但它可能把本来的不变张量也送到零。例如在特征 \(p\) 的域 \(k\) 上，一维空间 \(ke\) 的次数 \(p\) 张量幂中，全部置换作用都是恒等，故 \(N_p=p!\operatorname{id}=0\)，而张量幂、不变子空间和对称幂都非零。因此只删去分母不能补救平均公式；这里限制商映射反而是同构，失败的是这个求和算子。

\ProseParagraphBreak
正特征不仅影响平均过程，还使对称代数内部的取幂运算出现新的加法性。设 \(p\) 为素数，\(R\) 为特征 \(p\) 的交换含幺环，\(M\) 为幺性 \(R\) 模。对称代数交换，故
\[
(v+w)^p=v^p+w^p,\qquad (av)^p=a^pv^p.
\]
第一式由二项式系数 \(\binom pi\) 对 \(0<i<p\) 都被 \(p\) 整除得到。因此 \(\Phi_p:M\rightarrow\operatorname{Sym}_R^p(M)\)、\(v\mapsto v^p\) 是可加映射，并对 Frobenius 环同态 \(F_R:R\rightarrow R\)、\(a\mapsto a^p\) 半线性：\(\Phi_p(av)=F_R(a)\Phi_p(v)\)。这里半线性指映射可加，而标量 \(a\) 须先经过 \(F_R\) 才作用于像；若 \(F_R\) 为恒等，就得到通常的 \(R\) 线性。
\symidx{\(F_R\)：特征 \(p\) 环上的 Frobenius 同态}

\ProseParagraphBreak
若进一步 \(R=k\) 为域且 \(M\ne0\)，对非零 \(v\in M\)，先在 \(kv\) 上定义 \(av\mapsto a\)，再由\cref{b2:lem:functional-extension}延拓为 \(\ell:M\to k\)，便有 \(\ell(v)=1\)。对称多线性映射 \((m_1,\ldots,m_p)\mapsto\prod_i\ell(m_i)\) 由\cref{b2:prop:symmetric-power-universal}诱导线性映射 \(\operatorname{Sym}_k^p(M)\to k\)，将 \(v^p\) 送到 \(1\)，故 \(v^p\ne0\)。若 \(\Phi_p\) 是 \(k\) 线性的，便有 \((a^p-a)v^p=0\) 对每个 \(a\in k\) 成立；域上的非零向量不能被非零标量零化，故 \(a^p=a\)。因此在这个域且模非零的范围内，\(\Phi_p\) 线性当且仅当 \(F_k\) 为恒等。一般环上不能由 \(\Phi_p\) 线性就这样消去 \(v^p\)。因子对称、取对称商和取 \(p\) 次幂是相互联系但不同的运算。

\ProseParagraphBreak
第六章把扩张标量写成 \(S\otimes_R-\)。现在的代数构造只用有限张量、直和与生成关系，改变底环后也能与扩张标量相容；不必为模选择一组基。

\begin{proposition}\label{b2:prop:tensor-symmetric-base-change}
设 \(R,S\) 是交换含幺环，\(\varphi:R\to S\) 是保幺环同态，\(M\) 是幺性 \(R\) 模，记 \(P=S\otimes_RM\)。存在自然的分次 \(S\) 代数同构
\[
\begin{aligned}
\alpha_T&:S\otimes_RT_R(M)\xrightarrow{\ \cong\ }T_S(P),\\
\alpha_{\mathrm{Sym}}&:S\otimes_R\operatorname{Sym}_R(M)
\xrightarrow{\ \cong\ }\operatorname{Sym}_S(P).
\end{aligned}
\]
左侧的乘法均由 \((s\otimes a)(t\otimes b)=st\otimes ab\) 给出，次数 \(d\) 的分量分别为 \(S\otimes_RM^{\otimes_Rd}\) 与 \(S\otimes_R\operatorname{Sym}_R^d(M)\)。零次的同构均为 \(S\otimes_RR\to S\)、\(s\otimes r\mapsto s\varphi(r)\)。对 \(d\ge1\)，其公式为
\[
\begin{aligned}
\alpha_T\bigl(s\otimes(m_1\otimes_R\cdots\otimes_Rm_d)\bigr)
&=(s\otimes_Rm_1)\otimes_S(1\otimes_Rm_2)
\otimes_S\cdots\otimes_S(1\otimes_Rm_d),\\
\alpha_{\mathrm{Sym}}\bigl(s\otimes(m_1\cdots m_d)\bigr)
&=(s\otimes_Rm_1)(1\otimes_Rm_2)\cdots(1\otimes_Rm_d).
\end{aligned}
\]
当 \(d=1\) 时，右侧只保留首因子。不要求 \(S\) 作为 \(R\) 模平坦。
\end{proposition}
\begin{proof}
先给左侧代数结构。对任意含幺 \(R\) 代数 \(A\)，因 \(R\) 的结构像在 \(A\) 中心，所述乘法保持两处 \(R\) 平衡关系。例如
\[
(s\varphi(r)\otimes a)(t\otimes b)
=st\varphi(r)\otimes ab
=st\otimes (ra)b
=(s\otimes ra)(t\otimes b),
\]
另一处由 \(arb=abr\) 同样核对。因此乘法良定义；结合律与单位元 \(1\otimes1\) 由 \(S,A\) 的对应等式继承，\(s\mapsto s\otimes1\) 的像在中心，故得到 \(S\) 代数。若 \(A\) 分次，\cref{b2:prop:tensor-direct-sums}给出 \(S\otimes_RA\cong\bigoplus_d(S\otimes_RA_d)\)，且上述乘法使次数相加。

令 \(\iota_R:M\to T_R(M)\)、\(\iota_S:P\to T_S(P)\) 为一次分量包含。把 \(T_S(P)\) 经 \(\varphi\) 看成 \(R\) 代数，映射 \(m\mapsto\iota_S(1\otimes m)\) 为 \(R\) 线性：平衡关系 \(1\otimes rm=\varphi(r)\otimes m\) 保证标量相容。\cref{b2:thm:tensor-algebra-universal}给出 \(R\) 代数同态 \(h:T_R(M)\to T_S(P)\)。公式
\[
\alpha_T(s\otimes a)=s\cdot h(a)
\]
对 \(R\) 平衡，因为 \(h(ra)=\varphi(r)h(a)\)，所以诱导 \(S\) 线性映射。\(S\) 标量在目标中心，故
\[
\alpha_T\bigl((s\otimes a)(t\otimes b)\bigr)
=st\cdot h(ab)
=(s\cdot h(a))(t\cdot h(b)),
\]
并且它保单位元。\(h\) 在一次生成元上的值给出题述纯张量公式；将标量 \(s\) 移入第一个 \(P\) 因子即得到所写形式。

反向，\(S\) 线性映射
\[
j:P\to S\otimes_RT_R(M),\qquad s\otimes m\longmapsto s\otimes\iota_R(m)
\]
由 \(\operatorname{id}_S\otimes\iota_R\) 给出，故良定义。再次用张量代数的泛性质，得到 \(S\) 代数同态 \(\beta_T:T_S(P)\to S\otimes_RT_R(M)\)。它在纯张量上的公式为
\[
\beta_T\bigl((s_1\otimes m_1)\otimes_S\cdots\otimes_S(s_d\otimes m_d)\bigr)
=(s_1\cdots s_d)\otimes(m_1\otimes_R\cdots\otimes_Rm_d),
\]
在零次部分则为 \(s\mapsto s\otimes1\)。两个复合都固定 \(S\) 标量；\(\alpha_T\beta_T\) 在每个一次生成元 \(s\otimes m\in P\) 上为恒等，\(\beta_T\alpha_T\) 在 \(s\otimes\iota_R(m)\) 上为恒等。这些元素连同 \(S\) 生成相应代数，故两复合在全代数上为恒等。两向映射把一次生成元送到一次分量，把标量送到零次分量，所以保持分次。

对称情形中，把 \(m\) 送到 \(\operatorname{Sym}_S(P)\) 的一次生成元 \(1\otimes m\)，由\cref{b2:thm:symmetric-algebra-universal}延拓为 \(R\) 代数同态 \(h_{\mathrm{Sym}}:\operatorname{Sym}_R(M)\to\operatorname{Sym}_S(P)\)。再定义 \(\alpha_{\mathrm{Sym}}(s\otimes a)=s\cdot h_{\mathrm{Sym}}(a)\)；与张量情形相同的平衡等式和乘法等式说明它是 \(S\) 代数同态，且给出题述公式。

反向，\(S\otimes_R\operatorname{Sym}_R(M)\) 是交换 \(S\) 代数。将 \(s\otimes m\in P\) 送到目标代数一次分量 \(S\otimes_R\operatorname{Sym}_R^1(M)\) 中的对应元素，得到 \(S\) 线性映射；再用对称代数的泛性质，得到 \(\beta_{\mathrm{Sym}}:\operatorname{Sym}_S(P)\to S\otimes_R\operatorname{Sym}_R(M)\)。具体地，
\[
\beta_{\mathrm{Sym}}\bigl((s_1\otimes m_1)\cdots(s_d\otimes m_d)\bigr)
=(s_1\cdots s_d)\otimes(m_1\cdots m_d).
\]
两个复合在 \(S\) 标量和一次生成元上都为恒等，故在全代数上为恒等；生成元的次数也得到保留，因而同构分次。这里所用的是两个交换代数的泛性质，并未要求张量积保持单射。

对任意 \(R\) 线性映射 \(f:M\to M'\)，先扩张标量再作用于一次生成元，与先作用于 \(m\) 再按上述公式映射，都得到 \(s\otimes f(m)\)；在 \(S\) 标量上也相同。生成性因此给出
\[
\begin{aligned}
\alpha_{T,M'}(\operatorname{id}_S\otimes T_R(f))
&=T_S(\operatorname{id}_S\otimes f)\alpha_{T,M},\\
\alpha_{\mathrm{Sym},M'}(\operatorname{id}_S\otimes\operatorname{Sym}_R(f))
&=\operatorname{Sym}_S(\operatorname{id}_S\otimes f)\alpha_{\mathrm{Sym},M}.
\end{aligned}
\]
这证明两个同构对 \(M\) 自然。若 \(S'\) 也是交换含幺 \(R\) 代数，并给定保幺 \(R\) 代数同态 \(\psi:S\to S'\)，两侧由 \(s\mapsto\psi(s)\)、\(s\otimes m\mapsto\psi(s)\otimes m\) 诱导的映射，也与这些同构相容：上述纯张量公式两种复合均只将系数 \(s\) 换成 \(\psi(s)\)，生成性使相容性在全代数上成立。
\end{proof}
